\subsection{Candidate Validation and Soundness}
\label{app:validation-proofs}
This appendix summarizes how TEPL validates a completed attribute program
and why acceptance implies correctness. The guarantee covers the core
operations and exact arithmetic specified in
Section~\ref{sec:acceptance-soundness}; it does not claim
completeness of validation or verification of the compiler implementation.

\paperparagraph{Validation procedure.}
\label{app:validation-procedure}
The inputs are a rewrite $L\Rightarrow R$, its written condition $C$, a
source-metadata feasibility program $E_q$, a completed attribute program
$f$, and candidate obligations $\mathcal G$. Starting from source legality
and successful evaluation of $C$, the validator establishes that $E_q$
evaluates to a Boolean. Each short-circuit check is proved defined under
the preceding successful checks. The predicate cannot justify its own
initial definedness.

The validator then assumes a true result from $E_q$ and checks the
following properties. Other construction choices, candidate obligations,
and RHS legality remain goals to prove, not additional assumptions:
\begin{enumerate}
\item \emph{Defined attributes.} Original attribute expressions use captured
metadata and have valid types, list accesses, and integer calculations.
All RHS fields are filled, and supplied fields are preserved. Definedness
is checked before simplification, so an unused arithmetic result cannot
hide an invalid access.
\item \emph{Legal replacement.} Every RHS selector is distinct and in range;
batch and contracting roles are disjoint, and paired extents agree.
Coordinate mappings preserve axis identities and order. All obligations
in $\mathcal G$ must also be proved.
\item \emph{Equivalent result.} Both graphs are expressed as indexed input
reads, products, and finite sums. Checked algebraic laws normalize these
expressions. Comparison requires equal coefficients, input-read factors
including multiplicity, and reduction domains, together with equal
ordered output shapes.
\end{enumerate}
The validator returns \code{Proved} only when all checks have derivations.
An unresolved obligation or exhausted work limit returns \code{Unknown},
not a counterexample. At runtime, false feasibility is a normal reason to
skip a match; the builder is never evaluated on that path.

\paperparagraph{Metadata operations.}
The calculus interprets axis lists as finite ordered sequences.
$\gath(A,q)$ selects entries at positions $q$; $\pos(A,G)$ finds the
positions of distinct members $G$ in a distinct list $A$. Ordered
complement removes selected positions while preserving survivor order.
For distinct lists $A,B$, $\mpo(A,B)=(u,v)$ enumerates their common entries
in $A$'s order: $u$ contains their positions in $A$, and $v$ contains the
corresponding positions in $B$. Thus $\gath(A,u)=\gath(B,v)$.
The supported metadata expressions combine captured values, literals,
finite lists, concatenation, length, ranges, these sequence operations,
permutation inverse, integer addition and subtraction, and conditionals.
Bounds, membership, and permutation premises justify partial operations.
Strict operations require defined arguments; conditionals and Boolean
operators follow their short-circuit evaluation order.

\paperparagraph{Core reasoning rules.}
The validator uses facts from source legality, successfully evaluated
written conditions, and previously checked feasibility predicates. It
applies Boolean consequences and equality substitution on defined
expressions. A true disjunction does not establish either disjunct alone.
Sequence reasoning uses length and concatenation laws, gather composition,
position lookup, permutation inversion, paired matching, and ordered
complement. Each law follows from comparing lengths and entries; for
example, both sides of gather composition select $A[q[h[j]]]$.
Set inclusion and disjointness use the Boolean membership meanings of
union, intersection, and difference. Set equality alone never establishes
ordered equality. Integer reasoning uses exact ground calculations and
interval bounds for addition, subtraction, and lengths. Bounds must fit
the declared integer range; unsupported arithmetic remains unresolved.

For values, the supported laws are distributivity, associativity and
commutativity of exact addition and multiplication, reordering finite
sums, and bijective renaming of bound coordinates. Moving a factor across
a sum requires that it be independent of the bound coordinates. Renaming
preserves extents and avoids variable capture. Comparison retains read
multiplicity and reduction domains, including empty domains.

\begin{proof}[Proof sketch for Theorem~\ref{thm:validator-soundness}]
Suppose validation returns \code{Proved}, and fix a legal source match
$m$ satisfying $C$. The metadata rules preserve definedness and truth:
sequence laws preserve entries and order, selector premises justify
accesses, and integer bounds establish representability. Induction on
these derivations first proves Boolean evaluation of $E_q$. In a sequence
of checks, a false result stops evaluation; a true result establishes the
premise for the next check. No unexecuted branch is assumed defined.

If $E_q$ is false, runtime skips the match. If it is true, the same metadata
argument under the additional true predicates proves that $f$ is defined
and its attributes satisfy RHS legality. Executable expressions preserve
operator meanings and evaluation order, so $f\Downarrow_m\beta$ for the
established attributes $\beta$.

Lowering preserves each operation's value and ordered shape by its
declared semantics and validated coordinate mappings. Normalization uses
distributivity, reordering of finite sums, and bijective renaming of bound
coordinates. Moving factors across sums requires proved independence;
renaming preserves extents and avoids variable capture. These steps
preserve values under exact arithmetic, including empty sums. Successful
term comparison therefore establishes value equality for every compatible
input. The separate ordered-shape checks complete
$\operatorname{Valid}(m,\beta)$ under the stated dtype-compatibility
assumption. Since $m$ was arbitrary, the theorem follows.

The result applies to the operations and reasoning rules stated above.
Additional implementation rules require their own justification.
\end{proof}
